\section{\termFloerRegular{}参数}
\label{sec:floer-regular-parameters}

本节说明所谓“\termGenericParameter{}”的含义；在这些参数下，我们所考虑的各个模空间都成为光滑流形。这又引出各种\termFloerOperator{}的定义。

我们考虑满足
\[
  \phi_H^1(o_N)\pitchfork\nu^*S
\]
的对子 $(H,S)$。对任意这样的对子，下列方程的解只有有限多个：
\[
  \dot z=X_H(z),
  \qquad z(0)\in o_N,
  \quad z(1)\in\nu^*S.
\]
\footnote{译者注：原书此处把随时间变化的哈密顿方程写作 $\dot z=X_H(z)$，
时间参数 $t$ 在记号中脱漏；本章此前均作 $X_H(t,z)$。译文照录原式，特此说明。}
换言之，它们就是 $\Omega_S$ 上 $\mathcal A_H$ 的临界点。对给定的抽象紧流形 $S_0$，用 $\operatorname{Emb}(S_0:N)$ 表示 $S_0$ 到 $N$ 中的嵌入所成的集合，并引入其子集
\begin{equation}
\begin{split}
  \operatorname{Emb}^{H}(S_0:N)
  :=\{S\mid {}&\nu^*S\pitchfork\phi_H^1(o_N),\ %
  S\subset N\text{ 是 }S_0\text{ 的一个嵌入的像}\}.
\end{split}
\label{eq:transverse-embedding-space}
\end{equation}
对这样的 $S\in\operatorname{Emb}^{H}(S_0:N)$，我们还考虑 $S_0$ 的一个给定嵌入的\termIsotopyClass{}，并将它记为
\[
  \operatorname{Iso}^{H}(S_0:N)
  \subset\operatorname{Emb}^{H}(S_0:N).
\]

由标准横截性定理可知，在 $C^\infty$ 拓扑下，$\operatorname{Emb}^{H}(S_0:N)$ 在 $\operatorname{Emb}(S_0:N)$ 中稠密。接下来，我们考虑方程~\eqref{eq:perturbed-cauchy-riemann-conormal} 的解空间的\termFredholmRegularity{}。

现在，与定理 10.4.1 类似地证明下述定理。边界条件改变时的这类横截性结果已在 (Oh96a) 中证明，读者可参见其中的证明。

\begin{theorem}
\label{thm:floer-regular-parameters}
设 $H\in\mathcal H(S)$，其中\footnote{译者注：原书公式 (12.4.17) 把此前使用的 $\phi_H^1(o_N)$ 写成 $\phi_H(o_M)$，随后也把嵌入空间的底流形写成 $M$；本节没有说明这一字母和记号变更。译文照录。}
\begin{equation}
  \mathcal H(S)
  :=\{H\in\mathcal H\mid\phi_H(o_M)\pitchfork\nu^*S\},
  \label{eq:regular-hamiltonians-for-submanifold}
\end{equation}
并且 $S\in\operatorname{Emb}^{H}(S_0:M)$。则存在稠密子集
\[
  \mathcal J_{H,S}\subset\mathcal J^c,
\]
使得方程~\eqref{eq:perturbed-cauchy-riemann-conormal} 的所有解都是正则的。\footnote{译者注：原书把 $\mathcal J_{H,S}$ 写成 $\mathcal J^c$ 的子集，但式~(12.2.10) 使用的是时间依赖族 $J\in\mathcal P(\mathcal J^c)$；译文照录原式。}
\end{theorem}

我们还需要这个定理的参数化版本。

\begin{exercise}
\label{ex:parametrized-floer-regularity}
表述并证明\cref{thm:floer-regular-parameters}的参数化版本。
\end{exercise}

若三元组 $(H,J,S)$ 满足 $J\in\mathcal J_{H,S}$ 且 $H\in\mathcal H(S)$，则称它是\emph{\termFloerRegular{}}的。
